\documentclass[11pt]{article}
\usepackage[margin=1in]{geometry}
\usepackage{amsmath,amssymb,amsthm,mathtools}
\usepackage[T1]{fontenc}
\usepackage{lmodern,microtype}
\usepackage[colorlinks=true,linkcolor=blue!45!black,urlcolor=blue!45!black,citecolor=blue!45!black]{hyperref}
\usepackage{xcolor}
\setlength{\emergencystretch}{2em}
\numberwithin{equation}{section}
\newtheorem{proposition}{Proposition}[section]
\newtheorem{lemma}[proposition]{Lemma}
\theoremstyle{remark}
\newtheorem{remark}[proposition]{Remark}
\newcommand{\dd}{\mathrm d}
\newcommand{\NN}{\mathsf N}
\newcommand{\Ric}{\operatorname{Ric}}
\newcommand{\Id}{\mathrm I}

\title{Wave metrics, classical copies, and heavenly perturbations\\
\large Derivation notes for the double-copy theory draft}
\author{}
\date{21 November 2026}
\begin{document}
\maketitle

These notes separate established constructions in the literature from direct
calculations for the wave family. The Kerr--Schild prescription, mixed Weyl
factorization, and heavenly potential formulation are established background.
The comparisons with the manuscript's descended connection, the explicit
adapted-frame factorization below, and the displayed perturbation recursion
are calculations within that background. No novelty claim is made for these
specializations. All fields are holomorphic and local. The word ASD follows
the wave manuscript's convention: curvature vanishes on alpha planes.
References may use the opposite orientation and call the same sector SD.

\section{Fixed backgrounds, the Kerr--Schild copy, and the descended field}

Let
\begin{equation}\label{dc:wave}
 g_H=2\dd u\dd v+2\dd z\dd w+H(u,z)\dd u^2,
 \qquad
 \eta_0=2\dd u\dd v+2\dd z\dd w,
 \qquad
 \NN=\begin{pmatrix}0&0\\1&0\end{pmatrix}.
\end{equation}
The symmetric-product convention is $g_{uv}=g_{vu}=1$.
Here $\eta_0$ is the auxiliary flat background defined by the chosen
normal-form coordinates on the same coordinate domain. The one-form
$k=\dd u$ is null and parallel for $\eta_0$; its raised vector is
$k^\sharp=\partial_v$ and $k^\sharp(H)=0$. Thus
$g_H=\eta_0+Hk\otimes k$ is in Kerr--Schild form.

The manuscript's target metric is
$\eta=2(\dd a\dd d-\dd b\dd c)$. The fixed relabeling
\begin{equation}\label{dc:relabel}
 U=-b,\qquad V=c,\qquad Z=a,\qquad W=d
\end{equation}
puts it in the form $2\dd U\dd V+2\dd Z\dd W$.
The background comparison map sends
$(u,v,z,w)$ to $(U,V,Z,W)=(u,v,z,w)$. This is an explicitly chosen
coordinate comparison, not the direction-dependent generalized chart in the
manuscript. Flat-side formulas below use $u,v,z,w$ as labels after this
comparison; reinstating capital letters gives their target versions.

\begin{proposition}[Kerr--Schild copy in the chosen background]
The scalar $H$ and matrix-valued potential
\begin{equation}\label{dc:KS}
 \mathcal A_{\mathrm{KS}}=H\,\dd u\,\NN
\end{equation}
solve the flat massless scalar and source-free ASD Yang--Mills equations,
respectively. Their field equations reduce to linear equations in this
ansatz.
\end{proposition}
\begin{proof}
The scalar wave operator is
\begin{equation}\label{dc:flatbox}
 \Box_0=2(\partial_u\partial_v+\partial_z\partial_w),
 \qquad \Box_0 H=0.
\end{equation}
All matrix coefficients lie in the abelian subalgebra $\mathbb C\NN$, so
\begin{equation}\label{dc:FKS}
 F_{\mathrm{KS}}=\dd\mathcal A_{\mathrm{KS}}
 =H_z\,\dd z\wedge\dd u\,\NN.
\end{equation}
For a flat alpha plane use the generators
\begin{equation}\label{dc:flat-alpha}
 \ell_0(\zeta)=\partial_u+\zeta\partial_z,
 \qquad
 \ell_1(\zeta)=-\partial_w+\zeta\partial_v.
\end{equation}
Equation~\eqref{dc:FKS} vanishes on their span for every direction,
including the second projective chart. It is therefore ASD.
Its closedness and the ASD condition imply the source-free Maxwell
equations. Equivalently, the only possibly nonzero divergences differentiate
$H_z$ in $v$ or $w$, and hence vanish. The Yang--Mills equations reduce to
these Maxwell equations because all color commutators vanish.
\end{proof}

This is the pp-wave Kerr--Schild prescription of
\cite[Sec.~5.1, Eqs.~(72)--(74)]{DC-MOW}. The gravitational profile and
gauge profile are both $H$; the tensor structure changes from $k\otimes k$
to $k\NN$. Coupling and constant color normalizations have been suppressed.
The scalar $H$ is also a zeroth-copy profile: if $c^a$ and
$\widetilde c^{\tilde a}$ are fixed color vectors, the biadjoint field
$\Phi^{a\tilde a}=c^a\widetilde c^{\tilde a}H$ has vanishing quadratic
color interaction by antisymmetry of the structure constants, and satisfies
its field equation by \eqref{dc:flatbox}.

\begin{proposition}[Derivative relation; a fixed-frame calculation]
After \eqref{dc:relabel}, the manuscript's descended connection satisfies
\begin{equation}\label{dc:derivative}
 A_H=\frac12 H_z\,\dd u\,\NN
     =\frac12\partial_z\mathcal A_{\mathrm{KS}},
 \qquad
 F_{A_H}=\frac12 H_{zz}\,\dd z\wedge\dd u\,\NN.
\end{equation}
Here $\partial_z$ differentiates coefficients in the fixed coordinate and
matrix frames.
\end{proposition}
\begin{proof}
The manuscript gives
$A_H=-H_z(-b,a)\dd b\,\NN/2$.
Substituting \eqref{dc:relabel} gives \eqref{dc:derivative}.
Taking an exterior derivative proves the curvature formula.
\end{proof}

Both potentials are ASD Maxwell embeddings, but they are different fields.
In this fixed Kerr--Schild presentation $A_H$ is the conventional single
copy of the profile $H_z/2$, not of $H$. Its relation to the original
gravitational metric includes a transverse integration. In the displayed
frame this integration recovers $H$ up to $f(u)$, removable from $g_H$ by
a shift of $v$. Differentiation of a potential is not an invariant
operation on arbitrary gauge-equivalence classes; \eqref{dc:derivative}
therefore makes no such assertion.

\section{An adapted-frame mixed Weyl factorization}

\subsection{Separating color, spacetime spinors, and frame comparison}

The matrix $\NN$ acts on the transferred rank-two bundle. To obtain a
Maxwell two-form, choose its displayed generator and write
\begin{equation}\label{dc:strip}
 F_{A_H}=\mathcal F\,\NN,
 \qquad
 \mathcal F=\kappa\,\dd u\wedge\dd z,
 \qquad
 \kappa=-\frac12H_{zz}.
\end{equation}
This is a choice of color normalization. The spinor associated with
$\mathcal F$ is obtained using the \emph{flat spacetime} spin coframe; it
is not obtained by interpreting the two indices of $\NN$ as spacetime
spinor indices.

Choose the curved and flat spin coframes, respectively,
\begin{equation}\label{dc:two-coframes}
 (\theta^{AA'})=
 \begin{pmatrix}\dd u&\dd z\\-\dd w&\dd v+\frac12H\dd u\end{pmatrix},
 \qquad
 (\vartheta^{AA'})=
 \begin{pmatrix}\dd u&\dd z\\-\dd w&\dd v\end{pmatrix}.
\end{equation}
Each represents its metric as twice the determinant of the displayed
matrix of one-forms. The flat coframe is compatible with the target
coframe in the manuscript: after \eqref{dc:relabel},
\begin{equation}
 \begin{pmatrix}\dd a&\dd b\\\dd c&\dd d\end{pmatrix}
 =\begin{pmatrix}\dd U&\dd Z\\-\dd W&\dd V\end{pmatrix}
 \begin{pmatrix}0&-1\\1&0\end{pmatrix}.
\end{equation}
This constant primed-frame change does not identify the transferred
internal bundle with the tangent spin bundle.

For the following component comparison, identify the \emph{unprimed}
spin frames in \eqref{dc:two-coframes} by matching their labels.
Let $o_A=(1,0)$ in each frame and set $\epsilon_{0'1'}=1$.
Fix the component normalizations by
\begin{align}
 \mathcal F_{AA'BB'}&=\epsilon_{A'B'}f_{AB},
 \label{dc:Maxwell-spinor-convention}\\
 R_{AA'BB'CC'DD'}&=
 \epsilon_{A'B'}\epsilon_{C'D'}\Psi_{ABCD}
 \quad\hbox{for the Ricci-flat ASD metric }g_H.
 \label{dc:Weyl-spinor-convention}
\end{align}
The coordinate curvature convention is the manuscript's
$R_{uzuz}=-H_{zz}/2$. These equations fix any sign or numerical
normalization needed below; other curvature conventions merely change
the constant normalization of one factor or of the scalar.

\begin{lemma}[Aligned Maxwell fields]
For every holomorphic $r(u,z)$, the two-form
$r\,\dd u\wedge\dd z$ is a source-free ASD Maxwell field on both
$\eta_0$ and $g_H$. Every such $r$ also satisfies the scalar wave
equation on both backgrounds.
\end{lemma}
\begin{proof}
The two-form is closed. It vanishes on the flat alpha planes
\eqref{dc:flat-alpha} and on the curved alpha planes generated by
\[
 \partial_u-\frac12H\partial_v+\zeta\partial_z,
 \qquad -\partial_w+\zeta\partial_v.
\]
It is therefore ASD on both metrics, so its closedness gives the Maxwell
equations on each. The determinant of $g_H$ is constant, and its inverse
differs from $\eta_0^{-1}$ only by
$g_H^{vv}=-H$. Consequently
\begin{equation}\label{dc:curvedbox}
 \Box_{g_H}=\Box_0-H\partial_v^2,
\end{equation}
which, like $\Box_0$, annihilates $r(u,z)$.
\end{proof}

\begin{proposition}[Mixed Weyl relation in the specified frames]
With \eqref{dc:Maxwell-spinor-convention}--\eqref{dc:Weyl-spinor-convention},
the color-stripped descended field and gravitational curvature are
\begin{equation}\label{dc:aligned-spinors}
 f_{AB}=\kappa o_Ao_B,
 \qquad
 \Psi_{ABCD}=\kappa o_Ao_Bo_Co_D.
\end{equation}
Let $\sigma_{AB}=o_Ao_B$ be the Maxwell spinor of
$\dd u\wedge\dd z$ and let $S_{\mathrm{mix}}=1$. Then
\begin{equation}\label{dc:mixed-Weyl}
 \Psi_{ABCD}
 =\frac{f_{(AB}\sigma_{CD)}}{S_{\mathrm{mix}}}.
\end{equation}
Both Maxwell factors and the scalar satisfy their field equations on
the chosen flat background. The same component fields also satisfy
those equations on $g_H$.
\end{proposition}
\begin{proof}
In either coframe, $\dd u\wedge\dd z$ occupies the component
$(00',01')$, so \eqref{dc:strip} gives $f_{00}=\kappa$ with the other
components zero. All coordinate Riemann components are generated by
$R_{uzuz}=\kappa$ and curvature symmetries. Changing to the curved
dual frame introduces only multiples of $\partial_v$, which annihilate
the curvature, so the sole independent Weyl component is
$\Psi_{0000}=\kappa$. This proves \eqref{dc:aligned-spinors}; its
symmetrized product gives \eqref{dc:mixed-Weyl}. The preceding lemma
establishes the differential equations, including for $r=1$.
\end{proof}

The metric is Petrov type N where $\kappa\ne0$ and type O where
$\kappa=0$. Equation~\eqref{dc:mixed-Weyl} remains regular at both kinds
of points. It is a mixed Weyl factorization in the established sense of
\cite[Sec.~5.2, Eq.~(110)]{DC-LMNO}, specialized by a direct calculation
to the descended field. It is not a gauge-independent identification of
the transferred bundle with flat tangent spinors: it uses color
stripping, background coordinates, and the expressly chosen spin-frame
comparison. The manuscript's direction-dependent tangent-spinor map is
a different object.

On an open set where $\kappa$ is nonzero there is also a same-copy
factorization
\begin{equation}\label{dc:same-Weyl}
 \Psi_{ABCD}=\frac{f_{(AB}f_{CD)}}{S_{\mathrm{same}}},
 \qquad S_{\mathrm{same}}=\kappa.
\end{equation}
The scalar is a wave solution by the lemma. The quotient in
\eqref{dc:same-Weyl} is not defined at zeros of $\kappa$, although its
simplified expression extends. The mixed relation needs no division
by a varying field. Neither Weyl scalar here is asserted to be the
Kerr--Schild zeroth-copy profile $H$. Weyl copies of pp-waves have
functional freedom; see \cite[Sec.~IV]{DC-GGMPP}. Thus this
factorization is an exact compatibility result in an aligned sector,
not by itself evidence for a canonical correspondence or an interacting
color--kinematics construction.

\section{Heavenly potentials and the three different prepotentials}

\subsection{Conventions and direct comparison with SDYM}

For a scalar potential $\Theta(u,v,z,w)$ put
\begin{align}
 g_\Theta&=2\dd u\dd v+2\dd z\dd w
 +\Theta_{zz}\dd u^2+\Theta_{vv}\dd w^2
 -2\Theta_{vz}\dd u\dd w,
 \label{dc:heavenlymetric}\\
 \mathcal E(\Theta)&=
 \Box_0\Theta-\Theta_{zz}\Theta_{vv}+(\Theta_{vz})^2.
 \label{dc:heavenlyeq}
\end{align}
The equation $\mathcal E(\Theta)=0$ is a second-heavenly potential
formulation of Ricci-flat ASD gravity in this gauge. These conventions
are obtained from \cite[Eqs.~(2.14)--(2.20)]{DC-BGS} by $X=z$, $Y=-w$.
Choose a local double primitive $\Theta_0(u,z)$ satisfying
\begin{equation}\label{dc:primitive}
 \Theta_{0,zz}=H.
\end{equation}
Then $g_{\Theta_0}=g_H$. The linear term and the quadratic determinant
term in \eqref{dc:heavenlyeq} vanish separately.

For clarity, the corresponding SDYM potential equation can also be
checked directly. Let $\phi$ be matrix-valued, with
\begin{equation}\label{dc:YMpotential}
 \mathcal A(\phi)=\phi_z\dd u-\phi_v\dd w,
 \qquad A_v=A_z=0.
\end{equation}
The curvature components obey
\begin{equation}
 F_{zv}=0,\qquad F_{uv}=F_{zw}=-\phi_{vz},\qquad
 F_{uw}=-\phi_{uv}-\phi_{zw}+[\phi_v,\phi_z].
\end{equation}
Evaluation on \eqref{dc:flat-alpha} shows that ASD is equivalent to
\begin{equation}\label{dc:SDYM}
 \Box_0\phi-2[\phi_v,\phi_z]=0,
\end{equation}
the convention of \cite[Eqs.~(2.4)--(2.8)]{DC-BGS} under the same
coordinate change.

Define the spacetime Poisson bracket
\begin{equation}\label{dc:kinbracket}
 \{f,g\}_{\mathrm{kin}}=f_vg_z-f_zg_v.
\end{equation}
Then the gravitational equation becomes
\begin{equation}
 \Box_0\Theta-\{\Theta_v,\Theta_z\}_{\mathrm{kin}}=0.
\end{equation}
Up to the fixed factor of two in \eqref{dc:SDYM}, the nonlinear color
commutator is replaced by a kinematic Poisson bracket. This is an
equation-level comparison; a map between individual solutions requires
its own data and is not supplied merely by replacing notation.

In particular, the following three gauge prepotentials all solve
\eqref{dc:SDYM} but produce distinct fields:
\begin{align}
 \phi_{\mathrm{pot}}&=\Theta_0\NN,
 &\mathcal A(\phi_{\mathrm{pot}})&=\Theta_{0,z}\dd u\,\NN,
 \label{dc:potential-copy}\\
 \phi_{\mathrm{KS}}&=\Theta_{0,z}\NN,
 &\mathcal A(\phi_{\mathrm{KS}})&=H\dd u\,\NN,
 \label{dc:KS-potential}\\
 \phi_{\mathrm{desc}}&=\frac12\Theta_{0,zz}\NN=\frac12H\NN,
 &\mathcal A(\phi_{\mathrm{desc}})&=\frac12H_z\dd u\,\NN=A_H.
 \label{dc:desc-potential}
\end{align}
The first uses the same coefficient function as the gravitational
heavenly potential. The second is the conventional Kerr--Schild copy.
The third is the connection descended in the manuscript. They differ
by transverse differentiation, and must not be conflated. Additions
$a(u)z+b(u)$ to the double primitive change neither $H$ nor the metric;
where they affect a displayed fixed-color potential, they contribute
only a locally exact multiple of $\dd u$.

\subsection{A systematic perturbation calculation}

Write
\begin{equation}
 Q(f,g)=f_{vz}g_{vz}
 -\frac12(f_{zz}g_{vv}+f_{vv}g_{zz}),
 \qquad \mathcal E(\Theta)=\Box_0\Theta+Q(\Theta,\Theta).
\end{equation}
This symmetric bilinear form gives the direct expansion
\begin{equation}\label{dc:one-perturbation}
 \mathcal E(\Theta_0+\varepsilon\chi)
 =\varepsilon\mathcal L_H\chi
 +\varepsilon^2\big((\chi_{vz})^2-\chi_{zz}\chi_{vv}\big),
 \qquad
 \mathcal L_H=\Box_0-H\partial_v^2=\Box_{g_H}.
\end{equation}
Indeed, $2Q(\Theta_0,\chi)=-H\chi_{vv}$ by
\eqref{dc:primitive}, and \eqref{dc:curvedbox} identifies the operator.
An arbitrary $\chi$ is not automatically a solution.

For a formal series
\begin{equation}
 \Theta=\Theta_0+\sum_{n\ge1}\varepsilon^n\chi_n,
\end{equation}
the first two equations are
\begin{equation}\label{dc:second-order}
 \mathcal L_H\chi_1=0,
 \qquad
 \mathcal L_H\chi_2=
 \chi_{1,zz}\chi_{1,vv}-(\chi_{1,vz})^2.
\end{equation}
More generally, equating coefficients gives
\begin{equation}\label{dc:all-orders}
 \mathcal L_H\chi_n
 =\sum_{j=1}^{n-1}
 \left(\chi_{j,zz}\chi_{n-j,vv}
       -\chi_{j,vz}\chi_{n-j,vz}\right),
 \qquad n\ge2.
\end{equation}
These are formal necessary and sufficient coefficient equations; no
claim about convergence or a choice of initial data is included.
Already the first equation couples the perturbation to the background.
The second tests the perturbation's quadratic self-interaction, which
is absent when every $\chi_n$ depends only on $u,z$.

For a momentum covector $k$ and $x=(u,v,z,w)$, the plane wave
$e^{ik\cdot x}$ satisfies
\begin{equation}
 \{e^{ik_1\cdot x},e^{ik_2\cdot x}\}_{\mathrm{kin}}
 =X(k_1,k_2)e^{i(k_1+k_2)\cdot x},
 \qquad
 X(k_1,k_2)=k_{1z}k_{2v}-k_{1v}k_{2z}.
\end{equation}
Modes of $\Theta_0(u,z)$ have $k_v=k_w=0$, so $X$ vanishes on every
pair of modes in this sector. Hence this family cannot distinguish
nonzero Yang--Mills cubic factors proportional to $f^{abc}X$ from
gravitational factors proportional to $X^2$; both vanish here.
Allowing suitable $v$-dependent perturbations turns on this test.
Such a perturbation remains within the heavenly gravity formulation
provided the recursion is solved. It does not establish that the
manuscript's stronger full-connection transport extends to it.

\begin{thebibliography}{9}

\bibitem{DC-MOW}
R.~Monteiro, D.~O'Connell and C.~D. White,
\emph{Black holes and the double copy},
JHEP \textbf{12} (2014), 056.
\href{https://doi.org/10.1007/JHEP12(2014)056}{doi:10.1007/JHEP12(2014)056};
\href{https://arxiv.org/abs/1410.0239}{arXiv:1410.0239}.
The pp-wave prescription used here is Sec.~5.1, Eqs.~(72)--(74).

\bibitem{DC-LMNO}
A.~Luna, R.~Monteiro, I.~Nicholson and D.~O'Connell,
\emph{Type D spacetimes and the Weyl double copy},
Class. Quantum Grav. \textbf{36} (2019), 065003.
\href{https://doi.org/10.1088/1361-6382/ab03e6}{doi:10.1088/1361-6382/ab03e6};
\href{https://arxiv.org/abs/1810.08183}{arXiv:1810.08183}.
Mixed Weyl factorization is defined in Sec.~5.2, Eq.~(110).

\bibitem{DC-GGMPP}
H.~Godazgar, M.~Godazgar, R.~Monteiro, D.~Peinador Veiga and C.~N. Pope,
\emph{The Weyl double copy for gravitational waves},
Phys. Rev. Lett. \textbf{126} (2021), 101103.
\href{https://doi.org/10.1103/PhysRevLett.126.101103}{doi:10.1103/PhysRevLett.126.101103};
\href{https://arxiv.org/abs/2010.02925}{arXiv:2010.02925}.
The functional freedom of type N and pp-wave copies is discussed in Sec.~IV.

\bibitem{DC-BGS}
G.~R. Brown, J.~Gowdy and B.~Spence,
\emph{Self-dual fields on self-dual backgrounds and the double copy},
Phys. Rev. D \textbf{109} (2024), 026009.
\href{https://doi.org/10.1103/PhysRevD.109.026009}{doi:10.1103/PhysRevD.109.026009};
\href{https://arxiv.org/abs/2307.11063}{arXiv:2307.11063}.
Coordinate and potential conventions are in Sec.~2, especially
Eqs.~(2.4)--(2.8) and (2.14)--(2.23).

\end{thebibliography}
\end{document}
